%%%%%%%%%%%%%%%%%%%%%%%%%%%%%%%%%%%%%%%%%
% Plik z wstępem do pracy
% Szablon pracy dyplomowej
% Wydział Informatyki 
% Zachodniopomorski Uniwersytet Technologiczny w Szczecinie
% autor Joanna Kołodziejczyk (jkolodziejczyk@zut.edu.pl)
% Bardzo wczesnym pierwowzorem szablonu był
% The Legrand Orange Book
% Version 5.0 (29/05/2025)
%
% Modifications to LOB assigned by %JK
%%%%%%%%%%%%%%%%%%%%%%%%%%%%%%%%%%%%%%%%%


\chapter*{Wstęp}

Współczesne tempo życia oraz rosnąca liczba obowiązków sprawiają, że efektywne zarządzanie gospodarstwem domowym staje się wyzwaniem, szczególnie w przypadku osób dzielących przestrzeń życiową, takich jak studenci, współlokatorzy czy rodziny. Codzienne zadania, takie jak sprzątanie, zakupy, płacenie rachunków czy podejmowanie wspólnych decyzji, często stają się źródłem nieporozumień i konfliktów, gdy brakuje jasnego podziału odpowiedzialności i sprawnej komunikacji.

Tradycyjne metody organizacji, takie jak kartki na lodówce, ustne ustalenia czy chaotyczne konwersacje w komunikatorach internetowych, okazują się niewystarczające w obliczu złożoności współczesnego życia. Brak centralnego systemu do śledzenia wydatków, terminów płatności czy postępów w realizacji zadań domowych prowadzi do zapominania o ważnych terminach, dublowania zakupów oraz poczucia niesprawiedliwości w podziale obowiązków.

Celem niniejszej pracy jest zaprojektowanie i zaimplementowanie kompleksowego systemu informatycznego wspomagającego zarządzanie gospodarstwem domowym. Aplikacja ma na celu ułatwienie komunikacji między domownikami, automatyzację procesów związanych z płatnościami i zakupami oraz zapewnienie przejrzystości w podziale obowiązków.

Głównym elementem systemu jest aplikacja mobilna, która umożliwia użytkownikom tworzenie wirtualnych domów, zapraszanie do nich współlokatorów oraz wspólne zarządzanie zasobami. Kluczowe funkcjonalności obejmują system zarządzania zadaniami z możliwością przypisywania ich do konkretnych osób, moduł finansowy do śledzenia rachunków i wydatków (wspierany technologią OCR do automatycznego odczytu danych z faktur), wspólne listy zakupów oraz system ankiet do podejmowania demokratycznych decyzji.

Praca składa się z sześciu rozdziałów. Rozdział pierwszy zawiera analizę problemu oraz przegląd istniejących rozwiązań. Rozdział drugi przedstawia specyfikację wymagań funkcjonalnych i niefunkcjonalnych oraz projekt systemu. Rozdział trzeci opisuje wybrane technologie wykorzystane do realizacji projektu, w tym język Go dla warstwy serwerowej oraz framework React Native dla aplikacji mobilnej. Rozdział czwarty poświęcony jest szczegółom implementacyjnym, w tym strukturze bazy danych, architekturze API oraz rozwiązaniom z zakresu bezpieczeństwa. Rozdział piąty prezentuje wyniki testów i badań nad wytworzonym oprogramowaniem. Całość kończy podsumowanie i wnioski z realizacji projektu.